% !TEX root = ../main.tex

\chapter{Completeness and Topology in Synthetic Domain Theory}
\label{chapter-topology}

% This chapter will investigate the completeness conditions involved in SDT based on the orthogonality introduced in \cite{reus_general_1997} and results obtained in previous chapters. We will see how those conditions characterise the domain and topology in SDT. 
% % In this chapter, we will not automatically postulate the axioms introduced for $\I$. Instead, we will see the relationship between these axioms and the completeness conditions of $\I$.

\section{Orthogonality}
The Phoa principle and its variants are all statements of the form that ``a function can be uniquely identified by its values on a subset of its domain''. Expressed in commutative diagrams, that is
% https://q.uiver.app/#q=WzAsMyxbMCwwLCJYIl0sWzEsMCwiWSJdLFswLDEsIloiXSxbMCwxLCJmIl0sWzAsMiwiaCIsMl0sWzEsMiwiXFxvdmVybGluZXtofSIsMCx7InN0eWxlIjp7ImJvZHkiOnsibmFtZSI6ImRhc2hlZCJ9fX1dXQ==
\[\begin{tikzcd}
    X & Y \\
    Z
    \arrow["f", from=1-1, to=1-2]
    \arrow["h"', from=1-1, to=2-1]
    \arrow["{\overline{h}}", dashed, from=1-2, to=2-1]
\end{tikzcd}\]
where $h$ uniquely factors through $f$ as $h=\overline{h}\circ f$.
Reus and Streicher gave a generalisation of statements of such kind in \cite{reus_general_1997}.
\begin{definition}
    For $f:X\to Y$ and $g:Z\to W$, if for any $h:X\to Z$ and $k:Y\to U$, there is a unique $\alpha:Y\to Z$ such that the following diagram commutes:
    % https://q.uiver.app/#q=WzAsNCxbMCwwLCJYIl0sWzEsMCwiWSJdLFswLDEsIloiXSxbMSwxLCJVIl0sWzAsMSwiZiJdLFswLDIsImgiLDJdLFsxLDIsIlxcYWxwaGEiLDEseyJzdHlsZSI6eyJib2R5Ijp7Im5hbWUiOiJkYXNoZWQifX19XSxbMiwzLCJnIiwyXSxbMSwzLCJrIl1d
    \[\begin{tikzcd}
        X & Y \\
        Z & U
        \arrow["f", from=1-1, to=1-2]
        \arrow["h"', from=1-1, to=2-1]
        \arrow["\alpha"{description}, dashed, from=1-2, to=2-1]
        \arrow["k", from=1-2, to=2-2]
        \arrow["g"', from=2-1, to=2-2]
    \end{tikzcd}\]
    then $f$ is said to be \textbf{orthogonal} to $g$, denoted $f\perp g$. If $U=\top$ is the final object, $f$ is said to be \textbf{orthogonal} to $Z$, denoted $f\perp Z$, and $Z$ is said to be \textbf{right-orthogonal} to $f$, or \textbf{$f$-orthogonal}.
\end{definition}
An important result on completeness without any prior knowledge of $f$ is provided in \cite{reus_general_1997}.
\begin{theorem}
    \label{thm-prod}
    If for any $a\in A$, $B(a)$ is $f$-orthogonal, then $\prod_{a:A}B(a)$ is $f$-orthogonal.
\end{theorem}

There are several completeness properties involved in the foundations of domains in SDT. This chapter investigates the conditions for the interval type and, more generally, the observational algebras to satisfy those properties. This chapter will also attempt to illustrate the connections between the concept of sobriomorphisms and these properties.

\section{Orders from Paths and Completeness Properties}
\label{sect-order}

The purpose of introducing the interval type in SDT is to identify the information order on a type $X$ with the paths on it. We now see what properties we need to make the paths a domain.
\begin{definition}
    Let $A$ be a type and $x,y\in A$. The \textbf{path type} between $x$ and $y$, written as $\hom_A(x,y)$, $\hom(x,y)$, or $x\leadsto y$ if without ambiguity, is defined as
    $$\hom_A(x,y):=\set{p\in A^\I:p(0)=x\wedge p(1)=y}$$
\end{definition}
\begin{lemma}
    For any type $A$, $(-\leadsto-):A\to A\to \U$ is reflexive.
\end{lemma}
\begin{proof}
    $$\id:=\lambda x.\lambda\_.x:\forall(x:A).x\leadsto x$$
\end{proof}
\begin{lemma}
    Take any type $A$, $B$. $(-\leadsto-)$ is congruent to all functions $f:A\to B$.
\end{lemma}
\begin{proof}
    $$\lambda x.\lambda y.\lambda(p:x\leadsto y).f\circ p:\forall(x:A).\forall(y:A).(x\leadsto y)\to(f(x)\leadsto f(y))$$
\end{proof}
\begin{lemma}
    Take any type $A$, $B$, and two functions $f,g:A\to B$. If $f(x)\leadsto g(x)$, $f\leadsto g$.
\end{lemma}
\begin{proof}
    $$\lambda(P:\forall x.f(x)\leadsto g(x)).\lambda i.\lambda x.P(x)(i):(\forall x.f(x)\leadsto g(x))\to(f\leadsto g)$$
\end{proof}

So far we can see that the path type illustrates a reflexive relation that is persevered by all functions and has function extensionality on function types. We now introduce the completeness properties that force a type to be a synthetic category and, moreover, a synthetic poset. The first completeness property is the Segal-completeness that illustrates the composition of paths.
\begin{definition}
    A type $A$ is called \textbf{Segal-complete} (or \textbf{Segal} for short) if it is right-orthogonal to $\Lambda_2\hookrightarrow\Delta^2$. That is, every function of type $\Lambda_2\to A$ extends uniquely to a function of type $\Delta^2\to A$.
\end{definition}
% \begin{lemma}
%     If $A$ is Segal, the paths on $A$ are composable.
% \end{lemma}
% \begin{proof}
%     The composition $p\cdot q$ can be constructed using the following map $\Delta^2\to A$ extended from $\Lambda_2\to A$:
%     % https://q.uiver.app/#q=WzAsMyxbMCwxLCIwXFxtYXBzdG8geCJdLFswLDAsIjFcXG1hcHN0byB5Il0sWzEsMCwiMlxcbWFwc3RvIHoiXSxbMCwxLCJwIiwwLHsic3R5bGUiOnsiYm9keSI6eyJuYW1lIjoic3F1aWdnbHkifX19XSxbMSwyLCJxIiwwLHsic3R5bGUiOnsiYm9keSI6eyJuYW1lIjoic3F1aWdnbHkifX19XSxbMCwyLCJwXFxjZG90IHEiLDIseyJzdHlsZSI6eyJib2R5Ijp7Im5hbWUiOiJzcXVpZ2dseSJ9fX1dXQ==
%     \[\begin{tikzcd}
%         {1\mapsto y} & {2\mapsto z} \\
%         {0\mapsto x}
%         \arrow["q", squiggly, from=1-1, to=1-2]
%         \arrow["p", squiggly, from=2-1, to=1-1]
%         \arrow["{p\cdot q}"', squiggly, from=2-1, to=1-2]
%     \end{tikzcd}\]
% \end{proof}
In a Segal type, paths can be composed as shown in the following diagram:
% https://q.uiver.app/#q=WzAsMyxbMCwxLCIwXFxtYXBzdG8geCJdLFswLDAsIjFcXG1hcHN0byB5Il0sWzEsMCwiMlxcbWFwc3RvIHoiXSxbMCwxLCJwIiwwLHsic3R5bGUiOnsiYm9keSI6eyJuYW1lIjoic3F1aWdnbHkifX19XSxbMSwyLCJxIiwwLHsic3R5bGUiOnsiYm9keSI6eyJuYW1lIjoic3F1aWdnbHkifX19XSxbMCwyLCJwXFxjZG90IHEiLDIseyJzdHlsZSI6eyJib2R5Ijp7Im5hbWUiOiJzcXVpZ2dseSJ9fX1dXQ==
\[\begin{tikzcd}
    {1\mapsto y} & {2\mapsto z} \\
    {0\mapsto x}
    \arrow["q", squiggly, from=1-1, to=1-2]
    \arrow["p", squiggly, from=2-1, to=1-1]
    \arrow["{p\cdot q}"', squiggly, from=2-1, to=1-2]
\end{tikzcd}\]

Next, we introduce the Rezk-completeness following the definitions in \cite{buchholtz_synthetic_2022,sterling_domains_2025}.
\begin{definition}
    Let $\E$ be the colimit of the following diagram:
    % https://q.uiver.app/#q=WzAsNyxbMSwwLCJcXEkiXSxbMywwLCJcXEkiXSxbNSwwLCJcXEkiXSxbMiwxLCJcXERlbHRhXzIiXSxbNCwxLCJcXERlbHRhXzIiXSxbMCwxLCJcXHRvcCJdLFs2LDEsIlxcdG9wIl0sWzAsMywiXFxsYW1iZGEgaS4oaSxpKSIsMV0sWzEsMywiXFxsYW1iZGEgaS4oaSwwKSIsMV0sWzEsNCwiXFxsYW1iZGEgaS4oMSxpKSIsMV0sWzIsNCwiXFxsYW1iZGEgaS4oaSxpKSIsMV0sWzAsNV0sWzIsNl1d
    \[\begin{tikzcd}
        & \I && \I && \I \\
        \top && {\Delta^2} && {\Delta^2} && \top
        \arrow[from=1-2, to=2-1]
        \arrow["{\lambda i.(i,i)}"{description}, from=1-2, to=2-3]
        \arrow["{\lambda i.(i,0)}"{description}, from=1-4, to=2-3]
        \arrow["{\lambda i.(1,i)}"{description}, from=1-4, to=2-5]
        \arrow["{\lambda i.(i,i)}"{description}, from=1-6, to=2-5]
        \arrow[from=1-6, to=2-7]
    \end{tikzcd}\]
    A type $A$ is called \textbf{Rezk-complete} (or \textbf{Rezk} for short) if it is right-orthogonal to $\E\to\top$.
\end{definition}
Functions of type $\E\to A$ classify all isomorphism of the Segal type $A$, and the Rezk-completeness asserts that all isomorphisms are equivalences \cite{buchholtz_synthetic_2022}. If $A$ is both Segal and Rezk, $A$ is said to be a \textbf{synthetic category}.

To make a category a poset, we additionally require that parallel paths to be equivalent, which is described by $\I$-separateness \cite{sterling_domains_2025}.
\begin{definition}
    Let $\I_\parallel$ be the following pushout:
    % https://q.uiver.app/#q=WzAsNCxbMCwwLCJcXHNldHswLDF9Il0sWzEsMCwiXFxJIl0sWzAsMSwiXFxJIl0sWzEsMSwiXFxJX1xccGFyYWxsZWwiXSxbMCwxLCIiLDAseyJzdHlsZSI6eyJ0YWlsIjp7Im5hbWUiOiJob29rIiwic2lkZSI6InRvcCJ9fX1dLFswLDIsIiIsMix7InN0eWxlIjp7InRhaWwiOnsibmFtZSI6Imhvb2siLCJzaWRlIjoidG9wIn19fV0sWzIsM10sWzEsM10sWzMsMCwiIiwxLHsic3R5bGUiOnsibmFtZSI6ImNvcm5lciJ9fV1d
    \[\begin{tikzcd}
        {\set{0,1}} & \I \\
        \I & {\I_\parallel}
        \arrow[hook, from=1-1, to=1-2]
        \arrow[hook, from=1-1, to=2-1]
        \arrow[from=1-2, to=2-2]
        \arrow[from=2-1, to=2-2]
        \arrow["\lrcorner"{anchor=center, pos=0.125, rotate=180}, draw=none, from=2-2, to=1-1]
    \end{tikzcd}\]
    A type $A$ is called \textbf{$\I$-separated} if it is right-orthogonal to $\I_\parallel\to\I$.
\end{definition}
Functions of type $\I_\parallel\to A$ classify all pairs of parallel paths, and $\I$-separateness guarantees that there exists at most one path between any two points, which truncates a category to a poset\footnote{The antisymmetry is guaranteed by Rezk-completeness.}. An $\I$-separated synthetic category is called a \textbf{synthetic poset}.

In SDT, we care about posets with suprema for all $\omega$-chains, which gives an additional completeness property.
\begin{definition}
    A type $A$ is called \textbf{chain-complete} (or \textbf{complete} for short) if it is right-orthogonal to the embedding between the initial and final $L$-(co)algebras $\w\hookrightarrow\wb$.
\end{definition}
A chain-complete synthetic poset is called a \textbf{synthetic predomain}. If it in addition contains a minimum, it is called a \textbf{synthetic domain}.

\section{Topology and Completeness}
We now investigate under what conditions is the interval type a synthetic predomain. Notice that Theorem \ref{thm-prod} implies that an orthogonality property holds for all observational algebras iff it holds for the interval type $\I$. In this section, we only postulate Axiom \ref{ax-prop}, \ref{ax-lattice}, and \ref{ax-l-alg}, as we will relate Axiom \ref{ax-interpole} with completeness properties and other statements.

We first show that Axiom \ref{ax-interpole} is enough to illustrate a poset.
\begin{theorem}
    \label{thm-poset}
    $\I$ is a synthetic poset iff Axiom \ref{ax-interpole} holds.
\end{theorem}
\begin{proof}
    If $\I$ is a synthetic poset, for any $i\sqsubseteq j$, there is a unique path between $i$ and $j$, which implies Axiom \ref{ax-interpole}. In reverse, with Axiom \ref{ax-interpole}, we have $(i\sqsubseteq j)\cong(i\leadsto j)$, which makes $\leadsto$ a partial order on $\I$.
\end{proof}
Axiom \ref{ax-interpole} does not tell us whether $\I$ is chain-complete, so we have to make it an axiom.
\begin{axiom}
    \label{ax-omega}
    $\I$ is chain-complete.
\end{axiom}
\begin{theorem}
    \label{thm-domain}
    $\I$ is a synthetic domain iff Axiom \ref{ax-interpole} and \ref{ax-omega} holds.
\end{theorem}
\begin{proof}
    Derived immediately from Theorem \ref{thm-poset} and the fact that $0$ is the minimum of $\I$.
\end{proof}

Recall that we use the observational algebra $\Op(X):=\I^X$ to describe the topology of a space $X$, yet this structure only provides finite unions and intersections, while the common definition for a topology requires the open sets to be closed under arbitrary unions. Now, with the condition that $\I$ is a synthetic domain, we have countable unions given by
$$\bigsqcup_{n\in\N}\chi_n=\sup_{n\in\N}\bigsqcup_{k=0}^n\chi_k$$
Although this still differs from having arbitrary unions, one can see this as illustrating the topology of a second-countable space.

Now, we provide an alternative way to describe the completeness properties using the concept of sobriomorphisms.
\begin{definition}
    If $S\subseteq X$ and the inclusion $S\hookrightarrow X$ is a sobriomorphism, A is called a \textbf{spanning subset} of $X$, denoted $S\unlhd X$.
\end{definition}
\begin{lemma}
    If $A\unlhd B\unlhd C$, $A\unlhd C$.
\end{lemma}
\begin{proof}
    Derived immediately from the composition of isomorphisms $\Op(A)\cong\Op(B)\cong\Op(C)$.
\end{proof}
\begin{lemma}
    \label{lemma-intermediate}
    If $A\unlhd C$ and $A\subseteq B\subseteq C$, $A\unlhd B\unlhd C$.
\end{lemma}
\begin{proof}
    The restriction $(-|_A):\Op(C)\to\Op(A)$ factors as $\Op(C)\to\Op(B)\to\Op(A)$. Since $(-|_A)$ is invertible, so is $(-|_B):\Op(C)\to\Op(B)$ and $(-|_A):\Op(B)\to\Op(A)$.
\end{proof}
Recall that from Axiom \ref{ax-interpole} we have shown Theorem \ref{thm-sobrio} and \ref{thm-sobrio-omega}, which in reverse can be seen as a condition implying Axiom \ref{ax-interpole} itself.
\begin{lemma}
    $\Lambda_\omega\unlhd\Delta^\omega$ implies that $\I$ is a synthetic poset.
\end{lemma}
We also have an analogous condition for Axiom \ref{ax-omega}.
\begin{lemma}
    If $\I$ is a synthetic poset, $\Delta^\omega\unlhd\Delta^\infty$ implies that $\I$ is a synthetic domain.
\end{lemma}
Now we combine these two conditions using Lemma \ref{lemma-intermediate}.
\begin{theorem}
    \label{thm-last}
    $\Lambda_\omega\unlhd\Delta^\infty$ implies that $\I$ is a synthetic domain.
\end{theorem}
We now propose a conjecture from the observation of Theorem \ref{thm-last} to conclude the investigations on sobriomorphisms so far:
\begin{conjecture}
    \label{hypo-last}
    A type $A$ is a synthetic predomain if it is right-orthogonal to $\Lambda_\omega\hookrightarrow\wb$.
\end{conjecture}
Although this conjecture may not be true, we still wish that the gap between orthogonality and completeness of a synthetic predomain is reachable within a few additional properties. More specifically, we expect that, under certain situations, the orthogonality to the embedding above implies the Segal and chain completeness, since the embedding is the composition of $\Lambda_\omega\hookrightarrow\w$ and $\w\hookrightarrow\wb$.
